\section{Results}

\subsection{Effect of consistency configuration}

All normal controls passed. Figure~\ref{fig:prediction-observation} compares
the adversarial observations with the configuration predictions. C1
(\texttt{local}, \texttt{w:1}, no causal session) returned counterexamples for
all four properties. C6 (\texttt{majority}, \texttt{majority}, causal session)
produced no decidable histories for its four properties: its RYW and MR reads
timed out; the MW and WFR histories had unresolved write effects. These C6
outcomes are therefore not PASS results.
The complete matrix is in Appendix~\ref{app:full-results}.

\maybefigure[fig:prediction-observation]{submission/figures/prediction-observation.pdf}{Adversarial observations by configuration and property. Each cell represents five histories. When no outcome is decidable, its label gives the count of UNAVAILABLE, INDETERMINATE, and PRECONDITION\_MISS histories.}
\FloatBarrier

In RYW and MR, C1 returned stale values from a lagging member, while the
corresponding C6 reads timed out. In WFR, C3 histories passed and C4 produced a
counterexample.

\subsection{Effect of failures and network partition}

The fault campaign contains \RQTwoHistoryCount\ histories across
\RQTwoFaultEpisodeCount\ shared fault episodes. Table~\ref{tab:fault-summary}
summarizes the outcomes, client-operation success rate, and latency for each
fault.

No violations were observed under F1 or F2. All observed violations occurred
under the F3 replication partition. Under F3, none of the 20 C6 RYW histories
was decidable: W1 failed to complete within the time limit, so the dependent
read was never issued.

Under F3, violations were observed for RYW and MW under C1, and for WFR under
C1 and C4. In the C4 WFR case, the isolated former primary first accepted a
setup write of version 1 with \texttt{w:1}, and the client's local read
subsequently returned version 1. The dependent majority write was then sent to
the newly elected primary on the other side of the partition. Its pre-image
contained only version 0, showing that the write was performed on a state older
than the version previously read by the client. Under C3 and C6, the initial
read did not return version 1, so the dependent write W2 was not issued.

\begin{table}[ht]
\centering
\small
\setlength{\tabcolsep}{4pt}
\caption{Fault outcomes, client-operation success, and all-attempt p50/p95 latency.}
\label{tab:fault-summary}
\begin{tabular}{@{}l r r l l l@{}}
\toprule
Fault & Histories & Episodes & P/V/U/I/PM & Operation success ($S$) & p50/p95 (ms) \\
\midrule
\RQTwoFaultSummaryRows
\bottomrule
\end{tabular}
\end{table}

\subsection{Evidence from recorded histories}

Three paired contrasts isolate the effects of causal sessions, read concern,
and write concern. With read and write concerns held fixed, the C5 read
returned a stale value, whereas the matched C6 read with a causal session
carried \texttt{afterClusterTime} and timed out. For read concern, the C8
local read returned version 1, while the matched C5 majority read returned
version 0. For write concern, the C3 \texttt{w:1} write was acknowledged but
was later absent from W2's pre-image, whereas the matched C6 majority write
timed out. Each contrast had eight matched control/fault pairs. Selected pair
IDs and their recorded operations are provided in
Appendix~\ref{app:representative-histories}.

\maybefigure[fig:representative-trace]{submission/figures/representative-trace.pdf}{Representative C6 RYW history during a replication-path partition: the majority-acknowledged write completed, then the dependent causal majority read timed out without returning a value.}